% %%%%%%%%%%%%%%%%%%%%%%%%%%%%%%%%%%%%%%%%%%%%%%%%%%%%%%%%%%%%%%%%%%%%%%%%%%%%
% Thesis Defence
% Section: Cache Partitioning
% %%%%%%%%%%%%%%%%%%%%%%%%%%%%%%%%%%%%%%%%%%%%%%%%%%%%%%%%%%%%%%%%%%%%%%%%%%%%

\section[Irrégularités]{Approche 2 : Partitionnement du cache}

% \begin{frame}
%     (Reintro apres red section) (cette section: aspect irrégulier)

%     grosse matrice dans le cache -> thrashing (anim?)

%     2 cas triviaux ou cache ok -> petite matrice ou bande horizontal (rappeler relation entre c et cache du spmv)

%     Pas en pratique

%     Matrices creuses, éléments isolés -> Idée naive: cache devient fully associatif !
    
%     Mais pas possible, et structure exploitable, de la matrice x kernel
% \end{frame}

\begin{frame}
    \frametitle{aaa}

    Figure, thrashing

    Pour éviter, et cache fonctionne bien, petite matrice, ou bande horizontale (montrer sur la figure)

    Pas en pratique

    Mais cas proche existe, bandée (montrer sur la figure)

    Si juste bandée, pas ok, change rien

    Si parfaitement bandée -> cache ok, moyennant une formule

    Qu'est ce que ca donne apres étude sur le mem traffic -> 2 ordres de grandeur

    Idée : isoler une bande diagonale de la matrice -> DRC
\end{frame}

\begin{frame}[t]
    \frametitle{Multiplication Matrice-Vecteur dans un cache générique :}
    \small

    \begin{minipage}{0.48\linewidth}
        \raggedright

        \pause[1]%
        Matrice quelconque :

        \point{Évincements répétés dans le cache}

        \arrow{Le cache n'est pas utile} %Rouge

        \bigskip

        \pause[2]%
        Petite matrice%
        \pause[3]%
        , bande horizontale :

        \pause[4]%
        \point{Tiens dans le cache, pas d'évincement}

        \arrow{Le cache est utile} %Vert

        \arrow{Cas peu intéressants} %Rouge

    \end{minipage}
    \hfill
    \begin{minipage}{0.48\linewidth}

        \pause[1]%
        \includegraphics<1>[width=\linewidth]{Figures/PDF/LIST.pdf}
        \includegraphics<2->[width=\linewidth]{Figures/PDF/LogoUGA.pdf}

    \end{minipage}

\end{frame}

\begin{frame}[t]
    \frametitle{Multiplication Matrice-Vecteur dans un cache générique :}
    \small

    \begin{minipage}{0.48\linewidth}
        \raggedright

        \pause[1]%
        Matrice bandée :%ici, 2 anims(?) pour montrer ce qu'est une matrice bandée

        \point{Évincements répétés dans le cache}

        \pause[2]%
        \arrow{Le cache n'est pas utile} %Rouge

        \arrow{Structure très courante} %Vert

        \bigskip

        \pause[3]%
        Matrice parfaitement bandée :

        \point{Nombre d'évincements limité}

        \arrow{Le cache est utile} %Vert

        \arrow{Structure peu courante} %Rouge

    \end{minipage}
    \hfill
    \begin{minipage}{0.48\linewidth}

        \pause[1]%
        \includegraphics<1>[width=\linewidth]{Figures/PDF/LIST.pdf}
        \includegraphics<2->[width=\linewidth]{Figures/PDF/LogoUGA.pdf}

    \end{minipage}

\end{frame}

% \begin{frame}
%     Montrer des matrices

%     discuter de leur structures à gros grain et grain fin

%     uniform bandée parfaitement bandée

%     (dcd)
% \end{frame}

\begin{frame}[t]
    \frametitle{Traffic mémoire théorique avec des matrices bandée}
    \small

    \begin{minipage}{0.38\linewidth}
        \raggedright
        
        % Qu'est ce que ca donne apres étude sur le mem traffic -> 2 ordres de grandeur

        Graphe : Traffic mémoire du vecteur de sortie $c$ VS densité en dehors de la bande, pour différentes tailles de cache

        \point{densité dans la bande fixe (10\textsuperscript{-2})}

        \point{..}

        \pause[2]%
        \begin{highlight}{Conclusion}{Red}
            On sélectionne une bande diagonale de la matrice et on la met dans un cache set-associatif

            \point{}{<relation>}
        \end{highlight}

    \end{minipage}
    \hfill
    \begin{minipage}{0.58\linewidth}

        \pause[1]%
        <étiquette model C>

        \includegraphics<1>[width=\linewidth]{Figures/PDF/LIST.pdf}
        \includegraphics<2->[width=\linewidth]{Figures/PDF/LogoUGA.pdf}

        (graph avec courbes qui apparaissent au bon moment)

    \end{minipage}

\end{frame}

\begin{frame}
    Que fait on du reste ?

    Granularité

    Éléments isolés
    
    \point{Pas compatible avec ligne de cache}
    
    \arrow{Il faut cache par éléments}

    Pas de structure, hit imprévisible -> augmenter le hit rate dans espace de cache modeste : complètement associatif

    Transaction par éléments

    le reste -> SRT fully associatif, taille modeste

    Sur graphe, montrer aussi qu'il faut des transactions mémoire adaptées
\end{frame}

\begin{frame}
    Conclusion

    Deux régions pour traiter le dense et le creux

    C'est ma contribution
\end{frame}